\documentclass[10pt,aspectratio=169]{beamer}
\usepackage{ucas}

\title{UCAS Unofficial Beamer Template}
\subtitle{A Subtitle}
\author{\href{mailto:linchenghao21@mails.ucas.ac.cn}{Jonathan Lin}
\\ 
School of Artificial Intelligence, UCAS
}

\date{June 18th, 2023}

\begin{document}

\maketitle

\section{Introduction}

\begin{frame}{Acknowledgement}
    Special thanks to Rui Song from Hong Kong Polytechnic University (PolyU). 
    % 
    This template is based on \href{https://github.com/wurahara/PolyU-Beamer-Slides}{PolyU Beamer Slides}. This modification has been authorized. 
    
    Official resources, including logos, motto, calligraphy, are referenced from \href{https://onestop.ucas.edu.cn/home/info/6b9e95dc-5785-4eee-b25f-1f884698cfc3}{UCAS Onestop Site}. We have received authorization from UCAS for usage of these resources. 
    % 
    CASIA logo is from \href{https://www.overleaf.com/latex/templates/ucas-casia-beamer-theme/stdydfhvkctw}{UCAS (CASIA) Beamer Theme}. 
    % 
    Background pictures are from the internet. 

    \textbf{Not for commercial use!!!}

    Some useful info: 
    \begin{itemize}
        \item UCAS Blue: RGB(23, 73, 148)
        \item The background pictures are Yanqi Lake Campus of UCAS
    \end{itemize}
\end{frame}

\begin{frame}{Beamer for SINTEF slides}{\thesection \, \secname}
\begin{itemize}
\item We assume you can use \LaTeX; if you cannot,
\href{http://en.wikibooks.org/wiki/LaTeX/}{you can learn it here}
\item Beamer is one of the most popular and powerful document
classes for presentations in \LaTeX
\item Beamer has also a detailed
\href{http://www.ctan.org/tex-archive/macros/latex/contrib/beamer/doc/beameruserguide.pdf}{user manual}
\item Here we will present only the most basic features to get you up to speed
\end{itemize}
\end{frame}

\begin{frame}{Beamer vs. PowerPoint}
Compared to PowerPoint, using \LaTeX\ is better because:
\begin{itemize}
\item It is not What-You-See-Is-What-You-Get, but
What-You-\emph{Mean}-Is-What-You-Get:
you write the content, the computer does the typesetting
\item Produces a \texttt{pdf}: no problems with fonts, formulas,
      program versions
\item Easier to keep consistent style, fonts, highlighting, etc.
\item Math typesetting in \TeX\ is the best (e.g.: $\mathbb{W}_p$ distance):
\begin{equation}
    \mathbb{W}_{p}(\mathbf{p}\|\mathbf{q}) = \inf_{\gamma}\left\{ \int_{\mathcal{X}\times\mathcal{Y}}
    {\mathrm{d}(\bm{x}_i,\bm{y}_j)^{p}\mathrm{d}\gamma(\bm{x}_i,\bm{y}_j)}
    \mathrel{\left. \vphantom{\sum_{}} \middle\vert \right.} 
    \gamma\in\Gamma(\mathbf{p},\mathbf{q})\right\}^{1/p}
\end{equation}
\end{itemize}

\end{frame}

\section{Editing}

\begin{frame}[fragile]{Selecting the Class}
After the last update to the graphic profile, the \texttt{sintef} theme for
Beamer has been updated into a full-fledged class.
To start working with \texttt{sintefbeamer}, start a \LaTeX\ document with the preamble:
\begin{block}{Minimum SINTEF Beamer Document}
\begin{lstlisting}[language=TeX]
\documentclass{sintefbeamer}
\begin{document}
\begin{frame}{Hello, world!}
\end {frame}
\end{document}
\end{lstlisting}
\end{block}
\end{frame}

\begin{frame}[fragile]{Title page}
To set a typical title page, you call some commands in the preamble:
\begin{block}{The Commands for the Title Page}
\begin{lstlisting}[language=TeX]
\title{Sample Title}
\subtitle{Sample subtitle}
\author{First Author, Second Author}
\date{Defaults to today's}
\end{lstlisting}
\end{block}
You can then write out the title page with \verb|\maketitle|.

You can set a different background image than the default one with the
\verb|\titlebackground| command, set before \verb|\maketitle|.

In the \texttt{backgrounds} folder, you can find a lot of standard backgrounds
for SINTEF presentation title pages.

\end{frame}

\begin{frame}[fragile]{Writing a Simple Slide}
\framesubtitle{It's really easy!}
\begin{itemize}[<+->]
\item A typical slide has bulleted lists
\item These can be uncovered in sequence
\end{itemize}
\begin{block}{Code for a Page with an Itemised List}<+->
\begin{lstlisting}[language=TeX]
\begin{frame}
  \frametitle{Writing a Simple Slide}
  \framesubtitle{It's really easy!}
  \begin{itemize}[<+->]
    \item A typical slide has bulleted lists
    \item These can be uncovered in sequence
  \end{itemize}
\end{frame}\end{lstlisting}
\end{block}
\end{frame}

\begin{frame}[fragile]{Adding images}
\begin{columns}
\begin{column}{0.7\textwidth}
Adding images works like in normal \LaTeX:
\begin{block}{Code for Adding Images}
\begin{lstlisting}[language=TeX]
\usepackage{graphicx}
% ...
\includegraphics
[width=\textwidth]{images/default}
\end{lstlisting}
\end{block}
\end{column}
\begin{column}{0.3\textwidth}
\includegraphics[width=\textwidth]{source/ucas_logo_caligraphy.png}\\
\end{column}
\end{columns}
\end{frame}

\begin{frame}[fragile]{Splitting in Columns}
Splitting the page is easy and common;
typically, one side has a picture and the other text:
\begin{columns}
\begin{column}{0.6\textwidth}
This is the first column
\end{column}
\begin{column}{0.3\textwidth}
And this the second
\end{column}
\end{columns}
\begin{block}{Column Code}
\begin{lstlisting}[language=TeX]
\begin{columns}
    \begin{column}{0.6\textwidth}
        This is the first column
    \end{column}
    \begin{column}{0.3\textwidth}
        And this the second
    \end{column}
    % There could be more!
\end{columns}
\end{lstlisting}
\end{block}
\end{frame}

\begin{frame}[fragile]
\frametitle{Fonts}
\begin{itemize}
\item The paramount task of fonts is being readable
\item There are good ones...
  \begin{itemize}
  \item {\textrm{Use serif fonts only with high-definition projectors}}
  \item {\textsf{Use sans-serif fonts otherwise (or if you simply prefer them)}}
  \end{itemize}
\item ... and not so good ones:
  \begin{itemize}
  \item {\texttt{Never use monospace for normal text}}
\end{itemize}
\end{itemize}
\end{frame}

\begin{frame}[fragile]{Look}
\begin{itemize}
\item To change the colour of the title dash, give one of the class options
      \texttt{cyandash} (default), \texttt{greendash}, \texttt{magentadash},
      \texttt{yellowdash}, or \texttt{nodash}.
\item To change between the light and dark themes, give the class options
      \texttt{light} (default) or \texttt{dark}. It is not possible to switch
      theme for one slide because of the design of Beamer---and it's probably a
      good thing.
\item To insert a final slide, use \verb|\backmatter|.
\item The aspect ratio defaults to 16:9, but you can change it to 4:3 for old
      projectors by passing the class option \texttt{aspectratio=43}; any other
      values accepted by Beamer are also possible.
\end{itemize}
\end{frame}

\section{Summary}

\begin{frame}
\frametitle{Good Luck!}
\begin{itemize}
\item Enough for an introduction! You should know enough by now
\item If you have corrections or suggestions,
\href{mailto:linchenghao21@mails.ucas.ac.cn}{send them to me!}
\end{itemize}
\end{frame}

\begingroup
    \setbeamertemplate{headline}{}
    \setbeamertemplate{footline}{}
    \begin{frame}[c]{}
	\begin{figure}
            \centering
            \includegraphics[width=0.4\paperwidth]{source/ucas_lastpage.png}
        \end{figure}
        Motto of UCAS by former President Yongxinag Lu of CAS. \\
        I provide three translations of my own here (probably not accurate):
        \\
        English: \emph{Refined Wisdom, Diligent Aspiration, Unyielding Quest, Luminous Virtue}. 
        \\
        Français: \emph{Sagesse Raffinée, Aspiration Assidue, Quête Inflexible, Vertu Lumineuse}. 
        \\
        Latin: \emph{Multam Eruditionem, Firmam Voluntatem, Investigatio Rerum, Illuminatio Virtutis}.
    \end{frame}
\endgroup

\end{document}
